% ============================================================================
% LANGENTWURF LE-007: WLAN & Mobile Daten
% Profil: S (Senioren 65+) | Tier: 2 (Standard)
% ============================================================================

\documentclass[11pt,a4paper]{article}
\usepackage[utf8]{inputenc}
\usepackage[T1]{fontenc}
\usepackage[ngerman]{babel}
\usepackage{geometry}
\usepackage{fancyhdr}
\usepackage{lastpage}
\usepackage{xcolor}
\usepackage{tcolorbox}
\usepackage{enumitem}
\usepackage{longtable}
\usepackage{hyperref}
\usepackage{amssymb}

\geometry{left=2.5cm, right=2.5cm, top=2.5cm, bottom=2.5cm}
\definecolor{fach}{HTML}{b35c00}
\definecolor{gray}{HTML}{666666}
\definecolor{successgreen}{HTML}{2a7a4b}
\definecolor{warnred}{HTML}{c62828}

\tcbuselibrary{skins,breakable}
\newtcolorbox{scorebox}{ colback=white, colframe=fach, fonttitle=\bfseries, title=Didaktik-Score, breakable }
\newtcolorbox{warnbox}[1][]{ colback=warnred!5, colframe=warnred, fonttitle=\bfseries, title=#1, breakable }

\pagestyle{fancy}
\fancyhf{}
\fancyhead[L]{\small\textcolor{gray}{Digitale Grundlagen -- 007 WLAN \& Mobile Daten}}
\fancyhead[R]{\small\textcolor{gray}{LE Tier 2 / Profil S}}
\fancyfoot[C]{\small\thepage{} / \pageref{LastPage}}
\renewcommand{\headrulewidth}{0.4pt}

\begin{document}

\begin{titlepage}
    \centering
    \vspace*{2cm}
    {\Large\textcolor{gray}{Digitale Grundlagen -- Senioren 65+}}\\[2cm]
    {\Huge\textcolor{fach}{\textbf{Langentwurf}}}\\[0.5cm]
    {\large Tier 2 | Profil S}\\[2cm]
    {\LARGE\textbf{007 -- WLAN \& Mobile Daten}}\\[0.5cm]
    {\large Modul B: Verbindungen \& Sicherheit}\\[2cm]
    \begin{tabular}{rl}
        \textbf{Zeitrahmen:} & 60--75 Minuten \\
        \textbf{Vorkenntnisse:} & iPhone-Grundlagen, Kontrollzentrum \\
    \end{tabular}
    \vfill
    {\normalsize Stand: \today}
\end{titlepage}

\tableofcontents
\newpage

\section{Bedingungsanalyse}

\subsection{Ausgangssituation}
\begin{itemize}
    \item WLAN und mobile Daten werden oft verwechselt
    \item Angst vor hohen Kosten durch Datenverbrauch
    \item Unsicherheit: ``Bin ich gerade im WLAN?''
    \item Roaming im Ausland ist ein Kostenfalle-Risiko
\end{itemize}

\subsection{Typische Probleme}
\begin{itemize}
    \item ``Ich weiß nicht, ob ich WLAN habe''
    \item ``Mein Datenvolumen ist immer schnell aufgebraucht''
    \item ``Im Urlaub hatte ich eine riesige Rechnung''
    \item ``Was ist der Unterschied?''
\end{itemize}

\section{Sachanalyse}

\subsection{Kernkonzepte}
\begin{enumerate}
    \item \textbf{WLAN:} Funknetzwerk zu Hause/in Cafés -- meist kostenlos, begrenzte Reichweite
    \item \textbf{Mobile Daten:} Internet über Mobilfunknetz -- begrenzt (Datenvolumen), überall
    \item \textbf{Datenvolumen:} Monatliches Limit, danach gedrosselt oder Kosten
    \item \textbf{Roaming:} Mobilfunk im Ausland -- EU meist inklusive, außerhalb teuer!
    \item \textbf{Flugmodus:} Alles aus (Funkstille)
\end{enumerate}

\subsection{Alltagsanalogien}
\begin{tabular}{|l|p{8cm}|}
\hline
\textbf{Begriff} & \textbf{Analogie} \\
\hline
WLAN & Wasserhahn zu Hause (Festpreis, unbegrenzt) \\
Mobile Daten & Wasserflaschen unterwegs (begrenzte Anzahl) \\
Datenvolumen & Die Anzahl deiner Wasserflaschen pro Monat \\
Roaming & Wasser im Ausland kaufen (sehr teuer!) \\
Flugmodus & Alle Hähne zudrehen \\
\hline
\end{tabular}

\section{Didaktische Analyse}

\subsection{Stellung in der Reihe}
Erste Einheit in Modul B (Verbindungen \& Sicherheit). Grundlage für alle weiteren Online-Themen.

\subsection{Sicherheitsrelevanz}
\begin{warnbox}[Kostenfallen]
\begin{itemize}
    \item Roaming außerhalb der EU kann hunderte Euro kosten
    \item Automatische Downloads verbrauchen Datenvolumen
    \item Öffentliche WLANs können unsicher sein
\end{itemize}
\end{warnbox}

\section{Lernziele}

\begin{tabular}{|c|p{7cm}|c|c|}
\hline
\textbf{Nr.} & \textbf{Die Teilnehmer können...} & \textbf{Stufe} & \textbf{Niveau} \\
\hline
FZ1 & WLAN und mobile Daten unterscheiden & Verstehen & Basis \\
FZ2 & erkennen, ob sie gerade im WLAN sind & Anwenden & Basis \\
FZ3 & WLAN ein-/ausschalten und sich verbinden & Anwenden & Basis \\
FZ4 & ihr Datenvolumen prüfen & Anwenden & Basis \\
FZ5 & die Roaming-Gefahr erklären & Verstehen & Vertiefung \\
FZ6 & mobile Daten und Roaming deaktivieren & Anwenden & Vertiefung \\
\hline
\end{tabular}

\section{Verlaufsplanung}

\begin{longtable}{|p{1cm}|p{2cm}|p{5cm}|p{3cm}|p{2cm}|}
\hline
\textbf{Zeit} & \textbf{Phase} & \textbf{Inhalt} & \textbf{TN-Aktivität} & \textbf{Material} \\
\hline
0--10 & Einstieg & ``Was kostet Sie Ihr Internet?'' -- WLAN vs. Mobilfunk & Nachdenken & -- \\
\hline
10--25 & WLAN & Was ist WLAN? Wie verbinde ich mich? Statusleiste lesen & Mitmachen, verbinden & S.1 \\
\hline
25--40 & Mobile Daten & Was ist Datenvolumen? Wo sehe ich es? Kontrollzentrum & Prüfen am Gerät & S.1 \\
\hline
40--55 & Gefahren & Roaming erklären, Flugmodus, öffentliche WLANs & Zuhören & S.2 \\
\hline
55--70 & Übungen & Situationen durchspielen, Checkliste & Bearbeiten & S.3 \\
\hline
\end{longtable}

\section{Lösungen}

\textbf{Wo sehe ich, ob ich im WLAN bin?}
\begin{itemize}
    \item Statusleiste oben: WLAN-Symbol (Bogen) = WLAN aktiv
    \item Kontrollzentrum: WLAN-Symbol blau = verbunden
    \item Einstellungen $\rightarrow$ WLAN: Netzwerkname mit Häkchen
\end{itemize}

\textbf{Wie prüfe ich mein Datenvolumen?}
\begin{itemize}
    \item App des Mobilfunkanbieters (Telekom, Vodafone, O2)
    \item SMS an Anbieter (z.B. ``Verbrauch'' an Kurzwahl)
    \item Einstellungen $\rightarrow$ Mobilfunk $\rightarrow$ Aktueller Zeitraum
\end{itemize}

\section{Didaktik-Score}

\begin{scorebox}
\begin{tabular}{|c|l|c|c|p{3.5cm}|}
\hline
\# & Dimension & Gew. & Score & Notiz \\
\hline
1 & Differenzierung & 15\% & 8/10 & Basis/Vertiefung klar \\
2 & Taxonomie & 10\% & 9/10 & Verstehen + Anwenden \\
3 & Alltagsrelevanz & 10\% & 10/10 & Kostenfallen vermeiden! \\
4 & Aktivierung & 10\% & 9/10 & Am Gerät prüfen \\
5 & Scaffolding & 10\% & 9/10 & Wasser-Analogie hilft \\
6 & Feedback & 10\% & 8/10 & Situationen durchspielen \\
7 & Sprachliche Klarheit & 10\% & 9/10 & Keine Fachbegriffe ohne Erklärung \\
8 & Motivation & 10\% & 10/10 & Geld sparen motiviert! \\
9 & Barrierefreiheit & 10\% & 9/10 & Symbole erklärt \\
10 & Praktikabilität & 5\% & 9/10 & 70 Min angemessen \\
\hline
\multicolumn{3}{|r|}{\textbf{GESAMT:}} & \textbf{9.05/10} & \textcolor{successgreen}{\textbf{FREIGABE}} \\
\hline
\end{tabular}
\end{scorebox}

\end{document}
